%----------------------------------------------------------------------------------------
%   CLASS DEFINITIONS (EMBEDDED RESUME.CLS)
%----------------------------------------------------------------------------------------
\documentclass[9pt,letterpaper]{article} % Font size and paper type

\usepackage[parfill]{parskip} % Remove paragraph indentation
\usepackage{array} % Required for boldface (\bf and \bfseries) tabular columns
\usepackage{ifthen} % Required for ifthenelse statements
\usepackage[T1]{fontenc}
\usepackage{charter}
\usepackage{newunicodechar}
\newunicodechar{—}{---}
\newunicodechar{–}{--}
\usepackage{enumitem}
\usepackage{xcolor} % Recommended for more color options
\usepackage[colorlinks=true, urlcolor=blue]{hyperref}
\usepackage[left=0.5in,top=0.25in,right=0.4in,bottom=0.2in, columnsep=0.5cm]{geometry} % Tightened margins slightly
\usepackage{iftex}
\ifpdftex
  \input{glyphtounicode}
  \pdfgentounicode=1
\fi
\usepackage[none]{hyphenat}
\usepackage{multicol}
\pagestyle{empty} % Suppress page numbers

\makeatletter % Begin internal command redefinitions

%----------------------------------------------------------------------------------------
%	HEADINGS COMMANDS: Commands for printing name and address
%----------------------------------------------------------------------------------------

\def \name#1{\def\@name{#1}} % Defines the \name command to set name
\def \@name {} % Sets \@name to empty by default

\def \addressSep {$\diamond$} % Set default address separator to a diamond

% One, two or three address lines can be specified 
\let \@addressone \relax
\let \@addresstwo \relax
\let \@addressthree \relax

% \address command can be used to set the first, second, and third address (last 2 optional)
\def \address #1{
  \@ifundefined{@addresstwo}{
    \def \@addresstwo {#1}
  }{
  \@ifundefined{@addressthree}{
  \def \@addressthree {#1}
  }{
     \def \@addressone {#1}
  }}
}

% \printaddress is used to style an address line (given as input)
\def \printaddress #1{
  \begingroup
    \def \\ {\addressSep\ }
    \centerline{#1}
  \endgroup
  \par
}

% \printname is used to print the name as a page header
\def \printname {
  \begingroup
    \hfil{{\namesize\bf \@name}}\hfil
    \nameskip\break
  \endgroup
}

%----------------------------------------------------------------------------------------
%	PRINT THE HEADING LINES
%----------------------------------------------------------------------------------------

\let\ori@document=\document
\renewcommand{\document}{
  \ori@document  % Begin document
  \printname % Print the name specified with \name
  \@ifundefined{@addressone}{}{ % Print the first address if specified
    \printaddress{\@addressone}}
  \@ifundefined{@addresstwo}{}{ % Print the second address if specified
    \printaddress{\@addresstwo}}
     \@ifundefined{@addressthree}{}{ % Print the third address if specified
    \printaddress{\@addressthree}}
}

%----------------------------------------------------------------------------------------
%	SECTION FORMATTING
%----------------------------------------------------------------------------------------

% Defines the rSection environment for the large sections within the CV
\newenvironment{rSection}[1]{ % 1 input argument - section name
  \sectionskip
 \textbf{\Large #1} % Section title
  \sectionlineskip
  \hrule % Horizontal line
  \begin{list}{}{ % List for each individual item in the section
    \setlength{\leftmargin}{-0em} % Margin within the section
  }
  \item[]
}{
  \end{list}
}

%----------------------------------------------------------------------------------------
%	WORK EXPERIENCE FORMATTING
%----------------------------------------------------------------------------------------

\newenvironment{rSubsection}[4]{ % 4 input arguments - company name, year(s) employed, job title and location
 {\bf #1} \hfill {#2} % Bold company name and date on the right
 \ifthenelse{\equal{#3}{}}{}{ % If the third argument is not specified, don't print the job title and location line
  \\
  {\em #3} \hfill {\em #4} % Italic job title and location
  }
}{
  \par \vspace{0.1em}
}

% The below commands define the whitespace after certain things in the document - they can be \smallskip, \medskip or \bigskip
\def\namesize{\huge} % Size of the name at the top of the document
\def\addressskip{} % The space between the two address (or phone/email) lines
\def\sectionlineskip{\smallskip} % The space above the horizontal line for each section 
\def\nameskip{\vspace{-0.5em}} % The space after your name at the top
\def\sectionskip{} % The space after the heading section

\makeatother % End internal command redefinitions

%----------------------------------------------------------------------------------------
%   DOCUMENT START
%----------------------------------------------------------------------------------------
\name{Abhishek Nagaraja \vspace{.2em} \hrule} % Your name

\address{Dallas, Texas \textbullet\ work.abhishekn@gmail.com \textbullet\ {\href{https://www.linkedin.com/in/abhisheknagaraja/}{linkedin.com/in/abhisheknagaraja}} \textbullet\ {\href{https://www.abhisheknagaraja.com/}{www.abhisheknagaraja.com}}}

\begin{document}


%----------------------------------------------------------------------------------------
%   SUMMARY
%----------------------------------------------------------------------------------------
Technical product leader, AI systems builder, and forward-deployed engineer who architects multi-service platforms, ships agentic workflows, and embeds directly with customers and live deployments to turn ambiguous requirements into production systems—driving measurable outcomes across B2B SaaS, real-world event deployments, and automation. Experienced in LangGraph, RAG architectures, voice AI, product strategy, and stakeholder alignment.

%----------------------------------------------------------------------------------------
%   EDUCATION SECTION
%----------------------------------------------------------------------------------------
\begin{rSection}{Education \vspace{.1em}}
\begin{rSubsection}{University of Texas at Arlington}{Arlington, Texas, USA}{\textbf{Master of Science in Computer Science}}{August 2023 - May 2024}
{\\ \textbf{Master of Science in Engineering Management}} \hfill {May 2024 - May 2026}
\end{rSubsection}
\begin{rSubsection}{Jawaharlal Nehru Technological University}{Hyderabad, India}{\textbf{Bachelor of Technology in Computer Science and Engineering}}{August 2019 - July 2023}
\end{rSubsection}
\end{rSection}%----------------------------------------------------------------------------------------
%   TECHNICAL PROJECTS
%----------------------------------------------------------------------------------------
\begin{rSection}{Technical Products and Projects \vspace{.1em}}
\begin{rSubsection}{\href{https://logostaxsystems.com}{Logos Tax Systems: B2B SaaS AI Tax Platform}}{2025 - Present}{}{}
\begin{itemize}[left=-1.5em, labelsep=.1cm, labelwidth=.2cm, itemsep=0.0em, topsep=0pt]
\item Architected a \textbf{B2B SaaS} tax platform—a portable AI advisor + CPA marketplace, not a filing tool or CPA replacement—by designing a \textbf{9-repo} microservice architecture with a provider-agnostic \textbf{LLM abstraction layer}, on \textbf{GCP Cloud Run}.
\item Secured \textbf{3 production AI agents} on \textbf{GCP Cloud Run}—public, authenticated, and internal dev copilot—by enforcing \textbf{Supabase JWT} auth, org-based RBAC, and \textbf{248 RLS policies} across ~150 tables, grounded in an \textbf{82\%}-accurate pgvector RAG.
\item Reduced per-document AI cost to \textbf{~\$0.003} by designing a tiered document-parsing strategy—deterministic regex first, \textbf{Claude Haiku} vision fallback second—for W-2, 1099s, Schedule C/E, and Form 1040, with \textbf{PII sanitization} before every LLM call.
\item Codified \textbf{11 IRS tax-code strategies} (Augusta Rule, S-Corp, QBI, Bonus Depreciation, Cost Segregation) into automated calculators, by building a \textbf{Cross-Intelligence Service} that ranks synergistic \textbf{ROI} recommendations per client profile.
\item Integrated a multi-tenant tax-software platform with \textbf{OAuth2} (Intuit ProConnect) and CSV export (Drake Tax), by building CPA approval workflows across \textbf{5 providers} and \textbf{4 export formats} on \textbf{17 Supabase tables} with row-level security.
\end{itemize}
\end{rSubsection}
\begin{rSubsection}{\href{https://abhisheknagaraja.com/content/google-maps-teardown}{Google Maps: Product Strategy}}{}{}{}
\begin{itemize}[left=-1.5em, labelsep=.1cm, labelwidth=.2cm, itemsep=0.0em, topsep=0pt]
\item Translated qualitative data from \textbf{42 user interviews} into two core problem statements, revealing a critical user need for guaranteed \textbf{menu accuracy} by utilizing \textbf{UX journey mapping} to identify high-friction drop-off points.
\item Architected a \textbf{product strategy} to directly address validated user needs, by utilizing a roadmap for a \textbf{centralized menu system} with granular filtering to capture the high-value market segment with critical dietary requirements.
\end{itemize}
\end{rSubsection}
\begin{rSubsection}{\href{https://abhisheknagaraja.com/content/clash-of-clans-teardown}{Clash of Clans: Growth Strategy}}{}{}{}
\begin{itemize}[left=-1.5em, labelsep=.1cm, labelwidth=.2cm, itemsep=0.0em, topsep=0pt]
\item Proposed a \textbf{viral growth strategy} driving \textbf{5-10\% new organic user acquisition} by creating a native social sharing feature for key in-game content, by utilizing \textbf{referral loop optimization} and player behavioral data patterns.
\item Projected a \textbf{15-20\% increase in join rates} and \textbf{20-30\% increase in feature adoption} by designing a revamped \textbf{Clan Discovery engine} and transforming the feed by utilizing \textbf{data-driven insights}, metrics, and predictive UX patterns.
\end{itemize}
\end{rSubsection}
\end{rSection}%----------------------------------------------------------------------------------------
%   WORK EXPERIENCE SECTION
%----------------------------------------------------------------------------------------
\begin{rSection}{Work Experience}
\begin{rSubsection}{Program Manager Intern, University of Texas at Arlington - Texas, USA}{October 2023 - October 2025}{}{}
\begin{itemize}[left=-1.5em, labelsep=.1cm, labelwidth=.2cm, itemsep=0.0em, topsep=0pt]
\item Orchestrated the end-to-end \textbf{product lifecycle} for 'MavMarket', scaling to \textbf{200+ vendors} and \textbf{5,000+ attendees} while generating \textbf{\$120,000} in revenue, by executing a comprehensive \textbf{vendor acquisition strategy} and a 50+ page \textbf{AI-powered SOP} in Gemini.
\item Improved stakeholder \textbf{NPS by 35\%} and engagement by \textbf{25\%} by executing a high-fidelity feedback system for 100+ vendors and 200+ attendees, utilizing \textbf{user research} and data analytics to identify and resolve key operational friction points.
\item Founded the university's \textbf{Entrepreneurship Club} and co-founded \textbf{1 Million Cups Arlington} with City Govt and the local community, by building a startup network through 20+ workshops and connecting founders with local mentors and resources.
\end{itemize}
\end{rSubsection}
\begin{rSubsection}{Business Analyst Intern, NFC Solutions Private Limited - Hyderabad, India}{April 2022 - January 2023}{}{}
\begin{itemize}[left=-1.5em, labelsep=.1cm, labelwidth=.2cm, itemsep=0.0em, topsep=0pt]
\item Drafted \textbf{Product Requirement Documents (PRDs)} and coordinated task delegation for 4 key projects, serving as the bridge between stakeholders and technical teams by utilizing \textbf{Jira} to ensure consistent and on-time delivery cycles.
\item Improved development efficiency by 20\% by serving as Scrum Master, applying \textbf{Lean-Agile principles} to streamline workflows and ceremonies by utilizing \textbf{bi-weekly retrospectives} to identify, document, and remove complex process blockers.
\item Optimized the end-to-end \textbf{order lifecycle} for an e-commerce platform to improve customer satisfaction, by leveraging \textbf{automated inventory triggers} and pricing logic to reduce data processing errors and improve reliability by 15\%.
\item Demonstrated execution excellence as measured by zero-defect deployments, by proactively mitigating critical project risks and utilizing Gantt charts in Jira to manage 100\% on-time delivery across 4 complex enterprise projects.
\end{itemize}
\end{rSubsection}
\begin{rSubsection}{Founder and Managing Director, e-DAM Community - Hyderabad, India}{March 2021 - July 2023}{}{}
\begin{itemize}[left=-1.5em, labelsep=.1cm, labelwidth=.2cm, itemsep=0.0em, topsep=0pt]
\item Architected scalable solutions for community management, impacting over 15,000 users by leveraging \textbf{serverless architectures} and \textbf{automated data pipelines} to maintain high system availability and total data integrity at scale.
\item Directed \textbf{e-DAM}, \textbf{e-DAM Spotlight}, and \textbf{e-DAM CMS} by owning roadmaps, release governance, and stakeholder alignment, by utilizing \textbf{Agile methodologies}, sprint prioritization, and user-feedback loops to accelerate iteration across all three ventures.
\item Spearheaded community growth to 5,000+ members by performing \textbf{Market Analysis} and building a strategic network of industry partners, by utilizing \textbf{200+ technical workshops} to foster a thriving, scalable, and inclusive developer ecosystem.
\end{itemize}
\end{rSubsection}
\end{rSection}%----------------------------------------------------------------------------------------
%   OTHER PROJECTS
%----------------------------------------------------------------------------------------
\begin{rSection}{Other Projects \vspace{.1em}}
\begin{rSubsection}{\href{https://www.npmjs.com/package/@abhishek.nagaraja/catms}{CATMS (Coding Agents Team Management System): Open-Source NPM Package}}{}{}{}
\begin{itemize}[left=-1.5em, labelsep=.1cm, labelwidth=.2cm, itemsep=0.0em, topsep=0pt]
\item Authored and published an open-source NPM CLI (\textbf{@abhishek.nagaraja/catms}), scaling to \textbf{500+ downloads} by providing workflow scaffolding and a local \textbf{Model Context Protocol (MCP)} memory server across \textbf{Cursor}, \textbf{Claude}, and \textbf{Antigravity}.
\item Engineered a local RAG pipeline using \textbf{PostgreSQL}, \textbf{pgvector}, and \textbf{Ollama}, automatically chunking and embedding Markdown session logs from an Obsidian vault to provide context-aware retrieval.
\item Exposed dynamic agent tools (\texttt{cams\_query}, \texttt{cams\_ingest}, \texttt{cams\_backfill}), allowing agents to autonomously query historical decisions and backfill new context to solve long-term memory limits.
\end{itemize}
\end{rSubsection}

\begin{rSubsection}{\href{https://withthara.com/}{Thara: Multi-Agent AI Ecosystem}}{}{}{}
\begin{itemize}[left=-1.5em, labelsep=.1cm, labelwidth=.2cm, itemsep=0.0em, topsep=0pt]
\item Architected production \textbf{multi-agent orchestration} across \textbf{23 specialized agents} for tasks, calendar, journaling, and voice workflows, by engineering \textbf{LangGraph} on \textbf{FastAPI} with \textbf{PostgreSQL/pgvector} and \textbf{LlamaIndex}-backed semantic memory.
\item Expanded Thara across \textbf{8 client surfaces}—Telegram, iOS, Android, macOS, iPad, watchOS, React web, and voice—by shipping one \textbf{FastAPI/WebSocket} core with native adapters and \textbf{GitHub Actions CI/CD} deployed to \textbf{AWS ECS Fargate} with Docker.
\item Pioneered a \textbf{Pipecat + Twilio} voice stack with \textbf{Whisper STT}, \textbf{Kokoro TTS}, and \textbf{Silero VAD} on \textbf{FastAPI WebSockets}, by routing morning kickoff, overdue, habit-risk, and evening flows through the same task and calendar agents as mobile and desktop.
\end{itemize}
\end{rSubsection}
\end{rSection}%----------------------------------------------------------------------------------------
%   SKILLS SECTION
%----------------------------------------------------------------------------------------
\begin{rSection}{Skills and Tools}
\begin{itemize}[left=-1.5em, labelsep=.1cm, labelwidth=.2cm, itemsep=0.0em, topsep=0pt]
\item \textbf{Full-Stack \& AI Engineering:} Next.js, React, NestJS, FastAPI, Node.js, TypeScript, Python, PostgreSQL/pgvector, Supabase, LangGraph, RAG architectures, Twilio, GCP Cloud Run, AWS ECS Fargate, Vercel, Docker, GitHub Actions CI/CD.
\item \textbf{Data Platforms \& Reliability:} Snowflake, BigQuery, streaming ingestion patterns, data quality SLAs, anomaly detection, capacity planning for data services, FinOps awareness, lineage documentation, replayable pipelines, cost guardrails, batch versus streaming tradeoffs.
\item \textbf{Strategy, Finance \& GTM Analytics:} TAM/SAM/SOM modeling, unit economics, pricing elasticity, packaging strategy, competitive battlecards, sales enablement briefs, executive QBR storytelling, partnership term sheets, co-marketing governance, win-loss analysis, pipeline hygiene, discount governance, scenario planning.
\item \textbf{Program Operations \& Delivery:} Cross-functional program governance, written stakeholder communications, dependency mapping, CI/CD release coordination, software delivery sequencing, LLM evaluation guardrails, steering committee facilitation, executive readouts, quarterly business review preparation, design partnership rituals, launch readiness reviews, incident retrospectives, operational metrics scorecards.
\end{itemize}
\end{rSection}%----------------------------------------------------------------------------------------
%   CERTIFICATIONS SECTION
%----------------------------------------------------------------------------------------
\begin{rSection}{Certifications}
\begin{multicols}{2}
\begin{itemize}[left=-1.5em, labelsep=.1cm, labelwidth=.2cm, itemsep=0.0em, topsep=0pt]
\item Google: Agile Project Management
\item Google: Data Analytics Specialization
\item Google: Foundation of Project Management
\item Google: Project Initiation: Starting a Successful Project
\item Certified Scrum Product Owner (CSPO)
\item McKinsey.org Forward Program
\end{itemize}
\end{multicols}
\end{rSection}

\end{document}